% Part: first-order-logic
% Chapter: beyond
% Section: second-order-logic

\documentclass[../../../include/open-logic-section]{subfiles}

\begin{document}

\olfileid[tr]{fol}{byd}{sol}

\olsection{İkinci Dereceden Mantık}

İkinci dereceden mantığın dili, yalnızca bireylerden oluşan !!a{domain}
üzerinde değil, bu !!{domain} üzerindeki bağıntılar üzerinde de niceleme
yapmaya olanak verir. Birinci dereceden bir $\Lang{L}$ dili verildiğinde, her
$k$ için $R$ !!{variable}{}i eklenir; bunlar $k$-lı bağıntılar üzerinde değer
alır ve bu değişkenler üzerinde nicelemeye izin verilir. $R$, $k$-lı bir bağıntı için
!!a{variable} ve $t_1$, \dots,~$t_k$ sıradan (birinci dereceden) terimlerse,
$\Atom{R}{t_1,\dots,t_k}$ atomik !!{formula}{}dür. Bunun dışında !!{formula}s
kümesi, ikinci dereceden nicelemeye ilişkin ek hükümlerle birlikte, birinci
dereceden mantıktaki gibi tanımlanır. Yalnızca birinci dereceden terimler için
!!{identity} bulunduğuna dikkat edin: $R$ ile $S$, aynı $k$~aritesine sahip
bağıntı !!{variable}s{}iyse, $\eq[R][S]$ ifadesini aşağıdakinin kısaltması olarak
tanımlayabiliriz:
\[
\lforall[x_1 \dots][\lforall[x_k][(\Atom{R}{x_1, \dots, x_k} \liff
  \Atom{S}{x_1, \dots, x_k})]].
\]

İkinci dereceden mantığın kuralları, niceleyici kurallarını yalnızca yeni
ikinci dereceden değişkenleri kapsayacak biçimde genişletir. Ancak burada bu
değişkenlerin $\Lang{L}$'nin !!{predicate}s{}iyle ve daha genel olarak
$\Lang{L}$'nin !!{formula}s{}iyle nasıl etkileştiğini açıklarken biraz dikkatli
olmak gerekir. En azından bağıntı değişkenleri terim sayılır; dolayısıyla şu
biçimde çıkarımlar vardır:
\[
!A(R) \Proves \lexists[R][!A(R)]
\]
Fakat $\Lang{L}$, sabit bir $<$ bağıntı simgesi içeren aritmetik diliyse, şu
çıkarımın da geçerli olması beklenir:
\[
x < y \Proves \lexists[R][\Atom{R}{x,y}]
\]
ya da verilmiş bir $!A$ !!{formula}{}ü için,
\[
!A(x_1, \dots, x_k) \Proves \lexists[R][\Atom{R}{x_1,\dots,x_k}]
\]
Daha genel olarak şu biçimde çıkarımlara izin vermek isteyebiliriz:
\[
\Subst{!A}{\lambd[\vec x][!B(\vec x)]}{R} \Proves
\lexists[R][!A]
\]
Burada $\Subst{!A}{\lambd[\vec x][!B(\vec x)]}{R}$,
$\Atom{R}{t_1,\dots,t_k}$ biçimindeki her atomik !!{formula}{}ün $!A$ içinde
$!B(t_1, \dots, t_k)$ ile değiştirilmesinin sonucunu gösterir. Bu son kural,
bir {\em kavrama şeması}na, yani şu biçimde bir aksiyoma denktir:
\[
\lexists[R][\lforall[x_1, \dots, x_k][(!A(x_1, \dots, x_k) \liff
\Atom{R}{x_1, \dots, x_k})]],
\]
İkinci dereceden dildeki her $!A$ !!{formula}{}ü için böyle bir aksiyom vardır
ve burada $R$ serbest değişken değildir. (Alıştırma: $R$'nin $!A$ içinde
geçmesine izin verilirse bu şemanın tutarsız olduğunu gösterin!)

Mantıkçılar ``ikinci dereceden mantığın aksiyomları''ndan söz ederken genellikle
birinci dereceden mantığın ikinci dereceden niceleyici kurallarıyla asgari
genişletilmesini, kavrama şemasıyla birlikte kastederler. Ancak bu aksiyom ve
kuralların daha zayıf alt sistemlerini incelemek de çoğu zaman ilgi çekicidir.
Örneğin kavrama aksiyom şemasının tüm genelliğiyle \emph{impredikatif}
olduğuna dikkat edin: ikinci dereceden niceleyiciler içeren !!a{formula}
tarafından ``tanımlanan'' bir $\Atom{R}{x_1, \dots, x_k}$ bağıntısının varlığını
ileri sürmeye izin verir; bu niceleyiciler ise $R$'nin kendisini de içeren,
bu tür bütün bağıntılar kümesi üzerinde değer alır!{} Yirminci yüzyılın
başlarında Russell paradoksuna verilen yaygın bir tepki, suçu bu tür
tanımlara yüklemek ve matematiğin temellerini geliştirirken onlardan kaçınmak
oldu. $!A$ !!{formula}{}ünde ikinci dereceden niceleyicilerin kullanımı
yasaklanırsa, biraz daha zayıf olan {\em predikatif} bir kavrama biçimi elde
edilir.

Semantik bakımdan ikinci dereceden !!{structure}, dil için birinci dereceden
!!{structure} ile ikinci dereceden niceleyicilerin üzerinde değer aldığı
!!{domain} üzerindeki bir bağıntılar kümesinin birleşimi olarak düşünülebilir
(daha kesin söylersek, her $k$ için aritesi~$k$ olan bir bağıntılar kümesi
vardır). Elbette kavrama !!{derivation} sistemine dâhilse, ``ikinci dereceden
kısım''da kavrama aksiyomlarını sağlayacak kadar bağıntı bulunması da gerekir;
yoksa !!{derivation} sistemi sağlam değildir!{} Yeterince bağıntı bulunmasını
sağlamanın kolay bir yolu, ikinci dereceden kısmı birinci dereceden kısım
üzerindeki \emph{bütün} bağıntılardan oluşturmaktır. Böyle !!a{structure}
\emph{tam} olarak adlandırılır ve bir bakıma gerçekten dilin ``amaçlanan
!!{structure}''ıdır. Dikkatimizi tam !!{structure}s{}la sınırlandırırsak
\emph{tam} ikinci dereceden semantik denilen şeyi elde ederiz. Bu durumda bir
!!{structure} belirtmek, birinci dereceden kısmı belirtmeye indirgenir; çünkü
ikinci dereceden kısmın içeriği bundan örtük olarak çıkar.

Özetle, ikinci dereceden mantıktan söz ederken bir miktar belirsizlik vardır.
!!{derivation} sistemi bakımından şu ikisinden biri kastedilebilir:
\begin{enumerate}
\item Bazı kavrama aksiyomlarıyla birlikte ``asgari'' bir ikinci dereceden
  !!{derivation} sistemi.
\item Tam kavramayı içeren ``standart'' ikinci dereceden !!{derivation}
  sistemi.
\end{enumerate}
Semantik bakımından da şu ikisinden biri ilgi konusu olabilir:
\begin{enumerate}
\item !!a{structure}{}nın, birinci dereceden bir kısım ile kavrama aksiyomlarını
  sağlayacak kadar büyük ikinci dereceden bir kısımdan oluştuğu ``zayıf''
  semantik.
\item Yalnızca tam !!{structure}s{}ın ele alındığı ``standart'' ikinci dereceden
  semantik.
\end{enumerate}
Mantıkçılar hangi !!{derivation} sistemini ya da semantiği kastettiklerini
belirtmediklerinde, genellikle her listedeki ikinci maddeyi kastediyor olurlar.
Bu semantiği kullanmanın üstünlüğü, göreceğimiz üzere, birçok doğal matematiksel
yapının kategorik betimlemesini vermesidir; aynı zamanda !!{derivation} sistemi
oldukça güçlüdür ve bu semantiğe göre sağlamdır. Sakıncasıysa !!{derivation}
sisteminin bu semantiğe göre \emph{tam olmaması}dır; aslında etkin biçimde
verilmiş \emph{hiçbir} !!{derivation} sistemi tam ikinci dereceden semantiğe
göre tam değildir. Öte yandan, !!{derivation} sisteminin zayıflatılmış
semantiğe göre \emph{tam} olduğunu göreceğiz; bu, bir tümce kanıtlanabilir
değilse onun yanlış olduğu, tam olması gerekmeyen \emph{bir} !!{structure}
bulunduğu anlamına gelir.

İkinci dereceden mantığın dili oldukça zengindir. Tekli bağıntılar !!{domain}{}in
alt kümeleriyle özdeşleştirilebilir; böylece özellikle bu kümeler üzerinde
niceleme yapılabilir. Örneğin doğal sayılar için tümevarım tek bir aksiyomla
ifade edilebilir:
\[
\lforall[R][((\Atom{R}{\Obj{0}} \land \lforall[x][(\Atom{R}{x} \lif
    \Atom{R}{x'})]) \lif \lforall[x][\Atom{R}{x}])].
\]
Aritmetik dilinde $\Obj 0, \Obj \prime, +, \times$ ve $<$ simgelerinin
bulunduğunu varsayarsak, davranışlarını betimlemek için şu aksiyomları
ekleyebiliriz:
\begin{enumerate}
\item $\lforall[x][\lnot x' = \Obj 0]$
\item $\lforall[x][\lforall[y][(x' = y' \lif x = y)]]$
\item $\lforall[x][(x + \Obj 0) = x]$
\item $\lforall[x][\lforall[y][(x + y') = (x + y)']]$
\item $\lforall[x][(x \times \Obj 0) = \Obj 0]$
\item $\lforall[x][\lforall[y][(x \times y') = ((x \times y) + x)]]$
\item $\lforall[x][\lforall[y][(x < y \liff \lexists[z][y = (x + z')])]]$
\end{enumerate}
Tam ikinci dereceden semantiği kullandığımız sürece bu aksiyomların, yukarıdaki
tümevarım aksiyomuyla birlikte, aritmetiğin standart modeli olan
!!{structure}~$\Struct{N}$ için kategorik bir betimleme sağladığını göstermek
zor değildir. Bu aksiyomların doğru olduğu herhangi bir
!!{structure}~$\Struct{M}$ verilsin. Bir $f$ fonksiyonunu $\Nat$'tan
$\Struct{M}$'nin !!{domain}{}ine gidecek biçimde, $\Nat$ üzerinde olağan
özyinelemeyle tanımlayın; öyle ki $f(0) = \Assign{\Obj 0}{M}$ ve
$f(x+1) = \Assign{\prime}{M}(f(x))$ olsun. $\Nat$ üzerinde olağan tümevarımı ve
(1) ile~(2) aksiyomlarının
$\Struct M$ içinde geçerli olduğu olgusunu kullanarak $f$'nin !!{injective}
olduğunu görürüz. $f$'nin !!{surjective} olduğunu görmek için $P$,
$\Domain{M}$'nin $f$'nin görüntüsünde bulunan elemanları kümesi olsun.
$\Struct M$ tam olduğundan $P$ ikinci dereceden !!{domain} içinde yer alır.
$f$'nin kuruluşu
gereği $\Assign{\Obj 0}{M}$'nin $P$ içinde olduğunu ve $P$'nin
$\Assign{\prime}{M}$ altında kapalı olduğunu biliyoruz. Tümevarım aksiyomunun
$\Struct M$ içinde (özellikle $P$ için) geçerli olması, $P$'nin $\Struct M$'nin
birinci dereceden !!{domain}{}inin tamamına eşit olmasını güvence altına alır.
Bu, $f$'nin !!a{bijection} olduğunu gösterir. $f$'nin bir homomorfizma olduğunu
göstermek de $\Nat$ üzerinde olağan tümevarımı art arda kullanınca bundan daha
zor değildir.

Küme kuramsal bakımdan bir fonksiyon yalnızca özel bir bağıntı türüdür;
örneğin tekli bir $f$ fonksiyonu ikili bir $R$ bağıntısıyla
özdeşleştirilebilir; bu bağıntının $\lforall[x][\lexists![y][R(x,y)]]$
koşulunu sağlaması gerekir. Dolayısıyla fonksiyonlar
üzerinde de niceleme yapılabilir. Tam semantik kullanılarak sonsuz
!!{structure}s sınıfı, $\Struct M$ !!{structure}{}sı için $\Struct M$'nin
!!{domain}{}inden kendisinin öz bir alt kümesine bire bir bir fonksiyon bulunması
koşuluyla tanımlanabilir:
\[
\lexists[f][(\lforall[x][\lforall[y][(\eq[f(x)][f(y)] \lif \eq[x][y])]]
  \land \lexists[y][\lforall[x][\eq/[f(x)][y]]])].
\]
Bu tümcenin değillemesi de sonlu !!{structure}s sınıfını tanımlar.

Ayrıca doğrusal sıralama tanımına aşağıdakini ekleyerek iyi sıralamalar sınıfı
tanımlanabilir:
\[
\lforall[P][(\lexists[x][\Atom{P}{x}] \lif \lexists[x][(\Atom{P}{x}
    \land \lforall[y][(y < x \lif \lnot \Atom{P}{y})])])].
\]
Bu, ``küme'' ile ``tekli bağıntı'' özdeşleştirmesi altında, boş olmayan
her kümenin bir en küçük elemanı bulunduğunu bildirir. Başka bir örnek olarak,
köşelerin bağlantısız parçalara önemsiz olmayan hiçbir ayrımı bulunmadığını
söyleyerek çizgelerdeki bağlantılılık kavramı ifade edilebilir:
\[
\lnot \lexists[A][(\lexists[x][A(x)] \land \lexists[y][\lnot A(y)]
  \land \lforall[w][\lforall[z][((\Atom{A}{w} \land \lnot \Atom{A}{z})
      \lif \lnot \Atom{R}{w,z})]])].
\]
Bir başka örnek olarak, sonlu !!{structure}s sınıfının, !!{domain}{}i çift
büyüklükte olan alt sınıfını tanımlamayı alıştırma olarak deneyebilirsiniz.
Daha çarpıcısı,
rasyonel sayıları içeren tam bir sıralı cisim olarak gerçek sayıların kategorik
bir betimlemesi verilebilir.

Kısacası ikinci dereceden mantığın ifade gücü, birinci dereceden mantığınkinden
çok daha yüksektir. İyi haber bu; şimdi kötü habere geçelim. Tam ikinci dereceden semantiğe
göre tam olan etkin bir !!{derivation} sistemi bulunmadığını daha önce
belirtmiştik. İyi ya da kötü, birinci dereceden mantığın kompaktlık ve
L\"owenheim--Skolem teoremleri dâhil birçok özelliği burada yoktur.

Öte yandan tam ikinci dereceden semantikten vazgeçip daha zayıf semantiği
kabul etmeye hazırsak, asgari ikinci dereceden !!{derivation} sistemi bu semantiğe göre
tamdır. Başka bir deyişle, $\Proves$ simgesini ``asgari sistemde kanıtlar'' ve
$\Entails$ simgesini ``zayıf semantiğe göre mantıksal olarak gerektirir'' diye
okursak, $\Gamma \Entails !A$ olduğunda $\Gamma \Proves !A$ olduğunu
gösterebiliriz. !!{derivation} sistemine belirli kavrama aksiyomları dâhil
edilmek istenirse, semantik bu aksiyomları sağlayan ikinci dereceden yapılarla
sınırlandırılmalıdır: örneğin $\Delta$ bir kavrama aksiyomları kümesinden
(belki de tümünden) oluşuyorsa, $\Gamma \cup \Delta \Entails !A$ olduğunda
$\Gamma \cup \Delta \Proves !A$ elde edilir. Özellikle $!A$, ele aldığımız
kavrama aksiyomları kullanılarak kanıtlanamıyorsa, bu kavrama aksiyomlarının
yine de geçerli olduğu bir $\lnot !A$ modeli vardır.

Zayıf semantik için tamlık teoreminin geçerli olduğunu görmenin en kolay yolu,
ikinci dereceden mantığı şöyle çok sıralı bir mantık olarak düşünmektir. Bir
sıra, olağan ``birinci dereceden'' !!{domain} olarak yorumlanır; ardından her
$k$ için ``aritesi $k$ olan bağıntılar''dan oluşan !!a{domain} bulunur. Dilin,
yerleşik ``$\Atom{\Obj{true}_k}{R,x_1,\dots,x_k}$'' bağıntı simgelerini
içerdiğini varsayarız; bu simge $R$ bağıntısının $x_1$, \dots,~$x_k$ için geçerli
olduğunu savlamak üzere kullanılır. Burada $R$, ``$k$-lı bağıntı'' sırasından
bir değişken, $x_1$, \dots,~$x_k$ ise birinci dereceden sıranın nesneleridir.

Bu özdeşleştirmeyle zayıf ikinci dereceden semantik, özünde çok sıralı mantığın
olağan semantiğidir; çok sıralı mantığın birinci dereceden mantığa
gömülebileceğini de zaten gözlemledik. İki yöndeki çevirileri hesaba katarsak,
ikinci dereceden mantığın zayıf anlayışı gerçekte kılık değiştirmiş bir birinci
dereceden mantık biçimidir; burada !!{domain}, uygun aksiyomlarca yönetilen hem
``nesneleri'' hem de ``bağıntıları'' içerir.

\end{document}
